\section{Metodi}

\subsection{Dataset}

Il dataset utilizzato nel presente lavoro è disponibile pubblicamente
su Kaggle \cite{Kaggle} e deriva da registrazioni EEG effettuate
presso il Psychology and Psychiatry Research Center del Roozbeh
Hospital di Teheran.

Il campione comprende complessivamente 121 bambini, di cui 61 con
diagnosi di ADHD e 60 soggetti di controllo, con età compresa tra
7 e 12 anni. Per il gruppo di controllo sono stati esclusi soggetti con disturbi psichiatrici ed epilessia. Tutti i partecipanti erano destrimani.

La numerosità dei due gruppi è pressoché equivalente; il problema di
classificazione può quindi essere considerato sostanzialmente
bilanciato.


\subsubsection{Protocollo sperimentale}

Le registrazioni sono state effettuate durante un compito di
attenzione visiva. Ai bambini venivano mostrate immagini contenenti
personaggi e veniva chiesto di contarne il numero. Dopo la risposta
del partecipante veniva immediatamente presentata l'immagine
successiva.

La durata della registrazione non era prefissata, ma dipendeva dal
tempo impiegato da ciascun bambino per completare il compito. Per
questo motivo le registrazioni presentano durate differenti tra i
soggetti. Nel dataset originale, la durata media riportata è pari a

\begin{equation}
	154.59 \pm 58.13 \ \text{s}
\end{equation}

per il gruppo ADHD e a

\begin{equation}
	124.91 \pm 29.90 \ \text{s}
\end{equation}

per il gruppo di controllo.

Questa variabilità ha reso necessario, nella successiva fase di
preprocessing, selezionare per tutti i soggetti un intervallo
temporale della stessa durata.


\subsubsection{Acquisizione EEG}

I segnali EEG sono stati acquisiti con una frequenza di campionamento
pari a

\begin{equation}
	f_s = 128 \ \text{Hz},
\end{equation}
Il sistema comprendeva
19 elettrodi posizionati secondo il sistema internazionale 10--20.

I canali utilizzati nelle analisi sono:

\begin{center}
	\texttt{Fp1}, \texttt{Fp2}, \texttt{F3}, \texttt{F4},
	\texttt{C3}, \texttt{C4}, \texttt{P3}, \texttt{P4},
	\texttt{O1}, \texttt{O2}, \texttt{F7}, \texttt{F8},
	\texttt{T7}, \texttt{T8}, \texttt{P7}, \texttt{P8},
	\texttt{Fz}, \texttt{Cz}, \texttt{Pz}.
\end{center}

Il sistema 10--20 definisce le posizioni degli elettrodi sulla base
di distanze pari al 10\% o al 20\% di alcune misure anatomiche della
testa, considerate lungo le direzioni fronto-occipitale e
latero-laterale. Questa procedura consente di ottenere un
posizionamento riproducibile tra soggetti con dimensioni e forme
della testa differenti.

Le lettere presenti nei nomi dei canali indicano la regione dello
scalpo in cui l'elettrodo è collocato: \texttt{Fp} indica le
posizioni frontopolari, \texttt{F} quelle frontali, \texttt{C}
quelle centrali, \texttt{T} quelle temporali, \texttt{P} quelle
parietali e \texttt{O} quelle occipitali. Tale nomenclatura fornisce
un'indicazione topografica approssimativa e non deve essere
interpretata come una localizzazione precisa della sorgente corticale
che genera il segnale.

I numeri dispari identificano gli elettrodi situati sul lato sinistro
dello scalpo, mentre i numeri pari identificano quelli situati sul
lato destro. La lettera \texttt{z}, presente nei canali
\texttt{Fz}, \texttt{Cz} e \texttt{Pz}, indica invece le posizioni
collocate lungo la linea mediana.

Nel dataset originale gli elettrodi laterali erano indicati secondo
la precedente nomenclatura del sistema 10--20. Prima dell'analisi i
nomi sono stati uniformati alla nomenclatura attuale mediante le
seguenti corrispondenze:

\begin{equation}
	\texttt{T3}\rightarrow\texttt{T7},
	\qquad
	\texttt{T4}\rightarrow\texttt{T8},
	\qquad
	\texttt{T5}\rightarrow\texttt{P7},
	\qquad
	\texttt{T6}\rightarrow\texttt{P8}.
\end{equation}


\subsubsection{Struttura e limitazioni dei dati}

Il dataset è organizzato in formato tabellare. Ogni riga corrisponde
a un campione temporale, mentre le colonne contengono i valori
registrati dai 19 canali EEG, l'identificativo del soggetto e la
classe diagnostica.

Sebbene l'EEG sia stato acquisito durante un compito costituito da
più immagini successive, nel dataset non sono disponibili marcatori
temporali dettagliati relativi alla comparsa degli stimoli, alle
risposte dei partecipanti o alla correttezza delle singole risposte.

Per questa ragione le analisi sono state invece condotte sul segnale continuo, ricavando per ogni soggetto feature spettrali e di connettività che
descrivono complessivamente l'intervallo temporale considerato.

\subsection{Preprocessing del segnale}

Prima di estrarre le features è necessario eseguire una fase di preprocessing.  

La procedura comprendeva:

\begin{enumerate}
	\item uniformazione della durata delle registrazioni;
	\item conversione dei valori da microvolt a volt;
	\item assegnazione delle coordinate del montaggio standard 10--20;
	\item filtraggio passa-banda tra 1 e 45 Hz;
	\item applicazione del riferimento medio comune;
	\item decomposizione ICA e classificazione delle componenti.
\end{enumerate}


\subsubsection{Uniformazione della durata delle registrazioni}

Le registrazioni originali presentavano durate differenti, poiché il
tempo necessario per completare il compito dipendeva dalla velocità
di risposta del singolo partecipante.

L'utilizzo di segnali di durata diversa avrebbe potuto introdurre una
fonte di variabilità metodologica. Per garantire che tutti i
soggetti contribuissero alle analisi con la stessa quantità di
segnale, è stato quindi selezionato un intervallo temporale massimo comune.

Per ogni soggetto sono stati esclusi i primi 2 secondi della
registrazione e sono stati mantenuti i successivi 60 secondi. Il
segmento analizzato corrisponde pertanto all'intervallo:

\begin{equation}
	t \in [2,62)\ \mathrm{s}.
\end{equation}

L'esclusione dei primi 2 secondi è stata introdotta per ridurre la
possibile influenza della fase iniziale di assestamento della
registrazione e dell'avvio del compito. Non essendo disponibili
marcatori temporali dettagliati, tale esclusione non corrisponde alla
rimozione di uno specifico evento sperimentale, ma rappresenta una
scelta metodologica applicata in modo identico a tutti i soggetti.

Con una frequenza di campionamento pari a 128 Hz, ogni registrazione
preprocessata contiene quindi:

\begin{equation}
	N
	=
	60\ \mathrm{s}
	\cdot
	128\ \mathrm{campioni/s}
	=
	7680\ \mathrm{campioni}
\end{equation}

per ciascuno dei 19 canali.


\subsubsection{Conversione delle unità e assegnazione del montaggio}

Nel dataset i valori EEG erano espressi in microvolt. Poiché
MNE-Python rappresenta internamente i segnali EEG nelle unità del
Sistema Internazionale, i dati sono stati convertiti in volt:

\begin{equation}
	x_{\mathrm{V}}(t)
	=
	x_{\mathrm{\mu V}}(t)
	\cdot 10^{-6}.
\end{equation}

Per ciascun soggetto è stato successivamente costruito un oggetto
\texttt{RawArray} di MNE, organizzato nella forma
canali $\times$ campioni.

Ai canali sono state assegnate le coordinate del montaggio
\texttt{standard\_1020}. Le coordinate spaziali sono necessarie sia
per rappresentare le topografie EEG sia per fornire a ICLabel
informazioni sulla distribuzione delle componenti sullo scalpo.


\subsubsection{Filtraggio passa-banda}

Il segmento EEG selezionato è stato filtrato tra 1 e 45 Hz mediante
un filtro FIR (\textit{Finite Impulse Response}) progettato con il
metodo \texttt{firwin} e applicato a fase zero:

\begin{equation}
	1\ \mathrm{Hz}
	\leq
	f
	\leq
	45\ \mathrm{Hz}.
\end{equation}

Un filtro FIR calcola ciascun campione del segnale in uscita come una
combinazione lineare di un numero finito di campioni del segnale in
ingresso:

\begin{equation}
	y[n]
	=
	\sum_{k=0}^{L-1}
	h[k]x[n-k],
\end{equation}

dove $x[n]$ è il segnale originale, $y[n]$ il segnale filtrato,
$h[k]$ sono i coefficienti del filtro e $L$ rappresenta la lunghezza
del filtro.

La denominazione \textit{finite impulse response} deriva dal fatto
che la risposta del filtro a un impulso ha durata finita.

Il limite inferiore a 1 Hz attenua le derive molto lente della linea
di base, che possono derivare da variazioni dell'impedenza degli
elettrodi, movimenti lenti e instabilità strumentali. Queste
componenti possono presentare ampiezza elevata e influenzare sia la
stima della potenza spettrale sia la decomposizione ICA.

Il limite superiore a 45 Hz conserva le bande di frequenza utilizzate
nel progetto e attenua le componenti ad alta frequenza, parte
dell'attività muscolare e il rumore della rete elettrica localizzato
a 50 Hz. Il filtro non elimina necessariamente in modo assoluto tutte
le frequenze superiori a 45 Hz, ma ne riduce progressivamente
l'ampiezza.

La frequenza di campionamento del segnale è pari a 128 Hz. Secondo il
teorema del campionamento, la massima frequenza teoricamente
rappresentabile senza ambiguità è la frequenza di Nyquist:

\begin{equation}
	f_{\mathrm{Nyquist}}
	=
	\frac{f_s}{2}
	=
	\frac{128}{2}
	=
	64\ \mathrm{Hz}.
\end{equation}

Questo limite deriva dal fatto che sono necessari almeno due campioni
per rappresentare un ciclo di oscillazione. Componenti con frequenza
superiore a 64 Hz non possono essere rappresentate correttamente e
possono comparire artificialmente a frequenze inferiori attraverso
un fenomeno chiamato \textit{aliasing}.

La frequenza di taglio superiore di 45 Hz è quindi inferiore alla
frequenza di Nyquist e lascia un margine per la banda di transizione
del filtro.

I coefficienti di un filtro FIR possono essere scelti simmetrici,
ottenendo una risposta a fase lineare. In una normale applicazione
causale, questa proprietà introduce un ritardo temporale costante,
approssimativamente pari a:

\begin{equation}
	\tau_g
	=
	\frac{L-1}{2}
\end{equation}

campioni.

Nel presente lavoro il filtro è stato applicato a fase zero,
compensando tale ritardo e mantenendo il segnale filtrato temporalmente
allineato a quello originale. L'espressione ``fase zero'' non indica
che la fase fisiologica del segnale EEG venga annullata, ma che il
filtraggio non introduce uno spostamento di fase netto artificiale.


\subsubsection{Riferimento medio comune}

Dopo il filtraggio è stato applicato il riferimento medio. L'EEG non misura un potenziale elettrico assoluto, ma una differenza di potenziale rispetto a un riferimento. La scelta del riferimento influenza pertanto sia
l'ampiezza dei segnali sia le relazioni osservate tra i canali.

Per ogni istante temporale è stata calcolata la media dei 19 canali:

\begin{equation}
	\overline{x}(t)
	=
	\frac{1}{M}
	\sum_{j=1}^{M}x_j(t),
	\qquad M=19,
\end{equation}

e tale media è stata sottratta da ciascun canale:

\begin{equation}
	x_i^{\mathrm{CAR}}(t)
	=
	x_i(t)-\overline{x}(t).
\end{equation}

Il riferimento medio non elimina la necessità di un riferimento, ma
sostituisce quello originario con una stima distribuita sull'intero
insieme degli elettrodi. In questo modo si riduce la dipendenza del
segnale da una singola posizione di riferimento e si rende più
uniforme il confronto tra i canali.

\subsubsection{Independent Component Analysis}

\paragraph{Motivazione e principio generale}

Un elettrodo EEG non registra l'attività di una singola sorgente
cerebrale. Il segnale misurato in ogni posizione dello scalpo è una
combinazione di diversi contributi, tra cui attività neurale,
movimenti oculari, attività muscolare e rumore strumentale.

Una stessa sorgente può inoltre essere rilevata da più elettrodi. Un
battito di ciglia, per esempio, produce un segnale particolarmente
evidente nei canali frontali, ma il suo effetto può propagarsi anche
verso canali più posteriori.

L'Independent Component Analysis (ICA) è una tecnica di
\textit{blind source separation} che cerca di separare queste
sorgenti a partire dai soli segnali registrati, senza conoscere in
anticipo la loro forma.

Il modello ICA assume che i segnali osservati siano combinazioni
lineari di sorgenti non direttamente osservabili:

\begin{equation}
	\mathbf{X}
	=
	\mathbf{A}\mathbf{S},
\end{equation}

dove $\mathbf{X}$ contiene i segnali registrati dagli elettrodi,
$\mathbf{S}$ contiene le componenti indipendenti e $\mathbf{A}$ è
la matrice di mescolamento, che descrive come ogni componente si
proietta sui diversi canali.

L'algoritmo stima una matrice di separazione, indicata con
$\mathbf{W}$, attraverso la quale ricostruire le componenti:

\begin{equation}
	\widehat{\mathbf{S}}
	=
	\mathbf{W}\mathbf{X}.
\end{equation}


\paragraph{Indipendenza e non gaussianità}

L'ICA cerca componenti statisticamente il più possibile indipendenti.
L'indipendenza è una condizione più forte della semplice assenza di
correlazione: due segnali possono avere correlazione lineare nulla ma
presentare comunque una relazione statistica non lineare.

Nel contesto EEG, questo non significa che le sorgenti cerebrali siano
realmente indipendenti dal punto di vista fisiologico. L'indipendenza
viene utilizzata come criterio matematico per ottenere una
rappresentazione nella quale attività neurale e artefatti risultino
più separati rispetto ai segnali originali.

L'ICA sfrutta anche la non gaussianità delle sorgenti. In generale,
una combinazione di segnali indipendenti tende a essere più simile a
una distribuzione gaussiana rispetto ai segnali originari. Cercando
componenti con distribuzioni maggiormente non gaussiane, l'algoritmo
può quindi risalire alle sorgenti che hanno prodotto le miscele
registrate.

Gli artefatti oculari sono spesso adatti a questo tipo di separazione,
poiché i battiti di ciglia sono eventi relativamente brevi, di grande
ampiezza e con una distribuzione spaziale prevalentemente frontale.


\paragraph{PCA e numero di componenti}

Prima dell'ICA, MNE centra i dati, applica una procedura di whitening
ed esegue una Principal Component Analysis (PCA).

La PCA trasforma i dati in componenti non correlate e ordinate in base
alla quantità di varianza spiegata. L'ICA opera successivamente su
questo spazio cercando componenti non soltanto non correlate, ma il
più possibile indipendenti.


Nel presente lavoro è stato impostato:

\begin{equation}
	\texttt{n\_components}=0.95.
\end{equation}

Questo valore indica che, per ogni soggetto, viene selezionato il
numero minimo di componenti principali necessario a spiegare almeno
il 95\% della varianza dei dati. Il numero effettivo di componenti
può quindi variare leggermente tra i soggetti.

Con 19 canali e riferimento medio comune, il rango massimo teorico
dei dati è pari a 18, poiché la somma dei segnali riferiti alla media
è nulla a ogni istante.

Il restante contributo PCA non viene necessariamente eliminato dal
segnale: durante la ricostruzione finale MNE può reinserire la parte
dei dati non utilizzata per stimare l'ICA.


\paragraph{Ottimizzazione mediante Picard}

La decomposizione ICA è stata ottimizzata mediante Picard
(\textit{Preconditioned ICA for Real Data}).

Picard è un algoritmo numerico progettato per raggiungere la soluzione
ICA in modo rapido e stabile.

Nel presente lavoro è stata utilizzata la seguente configurazione:

\begin{verbatim}
	method = "picard"
	fit_params = {
		"ortho": False,
		"extended": True
	}
\end{verbatim}

L'opzione \texttt{extended=True} abilita una formulazione di tipo
Extended Infomax. Questa permette di trattare sia componenti
super-gaussiane, caratterizzate da eventi rari e di grande ampiezza,
sia componenti sub-gaussiane, caratterizzate da distribuzioni più
piatte.

L'opzione \texttt{ortho=False} consente a Picard di stimare una
matrice di separazione non vincolata all'ortogonalità, ottenendo una
configurazione compatibile con Extended Infomax.

Sono stati inoltre impostati:

\begin{equation}
	\texttt{random\_state}=42,
	\qquad
	\texttt{max\_iter}=3000.
\end{equation}

Il primo parametro rende riproducibile l'inizializzazione
dell'algoritmo, mentre il secondo definisce il numero massimo di
iterazioni consentite. La convergenza è stata verificata controllando
che l'algoritmo terminasse prima di raggiungere tale limite.


\paragraph{Classificazione delle componenti}

L'ICA produce una decomposizione statistica, ma non stabilisce
automaticamente quali componenti rappresentino attività cerebrale e
quali siano dovute ad artefatti.

Le componenti sono state quindi classificate mediante ICLabel, un
modello automatico che utilizza informazioni relative alla topografia,
allo spettro di potenza e all'andamento temporale di ciascuna
componente.

ICLabel assegna a ogni componente una probabilità di appartenenza a
categorie quali:

\begin{itemize}
	\item attività cerebrale;
	\item attività muscolare;
	\item artefatto oculare;
	\item attività cardiaca;
	\item rumore di linea;
	\item rumore di canale;
	\item altre sorgenti.
\end{itemize}

Il dataset non contiene canali EOG separati, che avrebbero permesso
di identificare direttamente gli artefatti oculari. La classificazione
automatica ha quindi utilizzato soprattutto la distribuzione spaziale,
le caratteristiche temporali e lo spettro delle componenti.

È stata adottata una procedura conservativa, sono state eliminate
soltanto le componenti classificate come oculari con probabilità
almeno pari a:

\begin{equation}
	p_{\mathrm{eye}}
	\geq
	0.90.
\end{equation}

Le componenti classificate come muscolari, rumore di linea o altre
sorgenti non sono state rimosse automaticamente, per ridurre il
rischio di eliminare anche attività cerebrale.


\paragraph{Ricostruzione del segnale}

Dopo aver identificato le componenti oculari, queste sono state
escluse dalla ricostruzione del segnale.

In termini schematici, le componenti selezionate vengono poste a zero
e il segnale viene ricostruito utilizzando tutte le componenti
rimanenti:

\begin{equation}
	\widehat{\mathbf{X}}_{\mathrm{clean}}
	=
	\widehat{\mathbf{A}}
	\widehat{\mathbf{S}}_{\mathrm{clean}}.
\end{equation}

Non vengono quindi eliminati interi elettrodi o interi intervalli
temporali. Da ogni canale viene sottratto soltanto il contributo delle
componenti identificate come artefatti oculari.


\paragraph{Limiti dell'ICA}

L'ICA fornisce una separazione statistica e non una ricostruzione
anatomica esatta delle sorgenti cerebrali. Una componente può
contenere contributi provenienti da più generatori, mentre una stessa
sorgente può essere distribuita tra componenti differenti.

La separazione è inoltre limitata dal numero di elettrodi disponibili.
Con 19 canali è possibile ottenere soltanto una rappresentazione a
bassa dimensionalità rispetto al numero reale di sorgenti neurali e
artefattuali.

L'ICA assume anche che la modalità con cui le sorgenti si proiettano
sugli elettrodi rimanga approssimativamente costante durante
l'intervallo analizzato. Movimenti degli elettrodi o variazioni
dell'impedenza possono ridurre la validità di questa ipotesi.

Infine, una componente classificata come artefattuale può contenere
anche una parte di attività neurale. Per questo motivo è stata
adottata una soglia elevata e una strategia conservativa di
esclusione.

\subsection{Estrazione delle feature spettrali}

L'estrazione delle feature spettrali è stata effettuata a partire
dai segnali EEG preprocessati. Per ciascun soggetto e per ciascuno
dei 19 canali, la serie temporale di 60 secondi è stata trasformata
dal dominio del tempo al dominio della frequenza mediante la stima
della densità spettrale di potenza
(\textit{Power Spectral Density}, PSD).

La rappresentazione temporale descrive come il potenziale elettrico
varia nel tempo, mentre la rappresentazione spettrale descrive come
la potenza del segnale è distribuita tra le diverse frequenze. Questo
passaggio è particolarmente utile per l'EEG, poiché numerosi fenomeni
neurofisiologici vengono descritti in termini di attività oscillatoria
all'interno di specifiche bande di frequenza.


\subsubsection{Stima della Power Spectral Density}

Per un segnale discreto $x[n]$, campionato alla frequenza $f_s$, la
trasformata discreta di Fourier permette di rappresentare il segnale
come combinazione di componenti sinusoidali:

\begin{equation}
	X[k]
	=
	\sum_{n=0}^{N-1}
	x[n]
	e^{-j2\pi kn/N}.
\end{equation}

Il modulo quadrato della trasformata fornisce una misura della potenza
associata alle diverse frequenze. Una stima basata direttamente
sull'intero segnale, detta periodogramma, presenta tuttavia una
varianza elevata: piccole variazioni del segnale possono produrre
oscillazioni considerevoli nella stima dello spettro.

Per ottenere una PSD più stabile è stato utilizzato il metodo di
Welch. Il segnale viene suddiviso in $K$ segmenti temporali
parzialmente sovrapposti. Per ciascun segmento viene calcolato un
periodogramma e le stime ottenute vengono successivamente mediate:

\begin{equation}
	\widehat{S}_{xx}(f)
	=
	\frac{1}{K}
	\sum_{k=1}^{K}
	P_k(f),
\end{equation}

dove $P_k(f)$ è il periodogramma del segmento $k$.

La media dei periodogrammi riduce la varianza della stima rispetto
all'utilizzo di un unico segmento. Il vantaggio viene ottenuto al
prezzo di una minore risoluzione temporale, che non costituisce però
un limite rilevante nel presente lavoro, poiché l'obiettivo è ottenere
una rappresentazione spettrale complessiva del soggetto e non seguire
l'evoluzione temporale delle bande durante il compito.

Sono stati utilizzati i seguenti parametri:

\begin{equation}
	n_{\mathrm{perseg}}=512,
	\qquad
	n_{\mathrm{overlap}}=256.
\end{equation}

Poiché la frequenza di campionamento è pari a 128 Hz, ogni segmento
contiene quattro secondi di segnale:

\begin{equation}
	T_{\mathrm{window}}
	=
	\frac{n_{\mathrm{perseg}}}{f_s}
	=
	\frac{512}{128}
	=
	4\ \mathrm{s}.
\end{equation}

L'overlap è pari al 50$\%$, poiché due finestre consecutive condividono
256 dei 512 campioni. La sovrapposizione permette di aumentare il
numero di periodogrammi disponibili per la media senza ridurre
ulteriormente la lunghezza delle finestre.

La distanza tra due frequenze consecutive della PSD è:

\begin{equation}
	\Delta f
	=
	\frac{f_s}{n_{\mathrm{perseg}}}
	=
	\frac{128}{512}
	=
	0.25\ \mathrm{Hz}.
\end{equation}

Una finestra di quattro secondi rappresenta quindi un compromesso tra
la necessità di distinguere le frequenze delle bande considerate e la
necessità di disporre di un numero sufficiente di segmenti da mediare.


\subsubsection{Calcolo delle potenze di banda}

A partire dalla PSD sono state calcolate le potenze assolute nelle
bande delta, theta, alpha e beta. Gli intervalli utilizzati
nell'estrazione delle feature spettrali sono riportati nella
Tabella~\ref{tab:eeg_bands}.

\begin{table}[H]
	\centering
	\begin{tabular}{lc}
		\hline
		\textbf{Banda EEG} & \textbf{Intervallo di frequenza} \\
		\hline
		Delta $(\delta)$ & 1--4 Hz \\
		Theta $(\theta)$ & 4--8 Hz \\
		Alpha $(\alpha)$ & 8--13 Hz \\
		Beta $(\beta)$ & 13--30 Hz \\
		\hline
	\end{tabular}
	\caption{Bande utilizzate per l'estrazione delle feature
		spettrali.}
	\label{tab:eeg_bands}
\end{table}

La potenza in una banda compresa tra $f_1$ e $f_2$ è definita come
l'area della PSD nell'intervallo corrispondente:

\begin{equation}
	P_{\mathrm{band}}
	=
	\int_{f_1}^{f_2}
	\widehat{S}_{xx}(f)\,df.
\end{equation}

Poiché la PSD è disponibile soltanto in corrispondenza di frequenze
discrete, l'integrale è stato approssimato numericamente mediante la
regola del trapezio:

\begin{equation}
	P_{\mathrm{band}}
	\approx
	\sum_{m}
	\frac{
		S(f_m)+S(f_{m+1})
	}{2}
	\left(f_{m+1}-f_m\right).
\end{equation}

Nel codice sono state incluse le frequenze che soddisfano:

\begin{equation}
	f\geq f_{\mathrm{min}},
	\qquad
	f<f_{\mathrm{max}}.
\end{equation}

Le potenze ottenute sono potenze assolute. 


\subsubsection{Theta/beta ratio}

Per ciascun canale è stato inoltre calcolato il rapporto tra la potenza
theta e la potenza beta:

\begin{equation}
	\mathrm{TBR}
	=
	\frac{P_{\theta}}{P_{\beta}}.
\end{equation}

Il rapporto è stato impostato come valore mancante nel caso,
puramente teorico, in cui la potenza beta fosse pari a zero. Nei dati
analizzati non sono stati osservati valori mancanti nelle feature
finali.


\subsubsection{Feature regionali}

Oltre alle feature dei singoli canali, sono state costruite feature
regionali aggregando gli elettrodi nelle seguenti cinque regioni:

\begin{itemize}
	\item frontale:
	\texttt{Fp1}, \texttt{Fp2}, \texttt{F3}, \texttt{F4},
	\texttt{F7}, \texttt{F8}, \texttt{Fz};
	
	\item centrale:
	\texttt{C3}, \texttt{C4}, \texttt{Cz};
	
	\item temporale:
	\texttt{T7}, \texttt{T8};
	
	\item parietale:
	\texttt{P3}, \texttt{P4}, \texttt{P7}, \texttt{P8},
	\texttt{Pz};
	
	\item occipitale:
	\texttt{O1}, \texttt{O2}.
\end{itemize}

Per ciascuna regione e banda, la potenza regionale è stata ottenuta
mediando le potenze dei canali appartenenti alla regione:

\begin{equation}
	P_{R,b}
	=
	\frac{1}{|R|}
	\sum_{c\in R}
	P_{c,b},
\end{equation}

dove $R$ indica la regione, $|R|$ il numero di canali appartenenti
alla regione e $b$ la banda di frequenza.

Il theta/beta ratio regionale è stato calcolato come rapporto tra la
potenza theta media della regione e la potenza beta media della stessa
regione:

\begin{equation}
	\mathrm{TBR}_{R}
	=
	\frac{P_{R,\theta}}{P_{R,\beta}}.
\end{equation}

Non è stata quindi calcolata la media dei TBR dei singoli canali;
prima sono state mediate separatamente le potenze theta e beta e
successivamente ne è stato calcolato il rapporto.


\subsubsection{Struttura del feature set spettrale}

Per ciascuno dei 19 canali sono state ottenute quattro potenze di banda
e un theta/beta ratio:

\begin{equation}
	19
	\times
	(4+1)
	=
	95
\end{equation}

feature a livello di canale.

Per ciascuna delle cinque regioni sono state ottenute quattro potenze
medie e un theta/beta ratio:

\begin{equation}
	5
	\times
	(4+1)
	=
	25
\end{equation}

feature regionali.

Il numero complessivo di predittori spettrali è quindi:

\begin{equation}
	95+25=120.
\end{equation}

La tabella finale contiene una riga per soggetto e 124 colonne:
120 feature predittive e quattro variabili descrittive,
\texttt{ID}, \texttt{Class}, \texttt{n\_samples} e
\texttt{duration\_sec}.


\subsection{Esplorazione delle feature spettrali}

Prima dell'addestramento dei classificatori è stata effettuata
un'analisi esplorativa con due obiettivi principali:

\begin{enumerate}
	\item descrivere la direzione e l'ampiezza delle differenze tra
	soggetti ADHD e controlli per ciascuna feature;
	
	\item valutare il grado di ridondanza tra le feature estratte.
\end{enumerate}


\subsubsection{Test di Mann--Whitney U}

Per confrontare i due gruppi è stato utilizzato il test di
Mann-Whitney U, un test non parametrico per due campioni indipendenti.

La scelta di un test non parametrico è motivata dal fatto che le
potenze EEG e i rapporti tra potenze possono presentare distribuzioni
asimmetriche e valori estremi. Non è stata quindi assunta la normalità
delle feature richiesta dal comune test $t$ parametrico.

Per ogni feature, i valori dei soggetti ADHD e dei controlli vengono
considerati congiuntamente, ordinati e trasformati in ranghi. Per il
primo gruppo, la statistica può essere espressa come:

\begin{equation}
	U_1
	=
	R_1
	-
	\frac{n_1(n_1+1)}{2},
\end{equation}

dove $R_1$ è la somma dei ranghi del primo gruppo e $n_1$ la sua
numerosità. Analogamente si ottiene:

\begin{equation}
	U_2
	=
	n_1n_2-U_1.
\end{equation}

Il test valuta se le osservazioni di uno dei due gruppi tendono a
occupare ranghi sistematicamente più alti o più bassi rispetto
all'altro gruppo.

È stato utilizzato un test a due code:

\begin{equation}
	H_0:
	\text{assenza di una differenza sistematica tra i gruppi},
\end{equation}

\begin{equation}
	H_1:
	\text{presenza di una differenza in una delle due direzioni}.
\end{equation}

Non è stata imposta a priori una direzione, poiché, nonostante le
ipotesi formulate sulla base della letteratura, il task attentivo
poteva produrre differenze diverse da quelle osservate nelle
registrazioni a riposo.

Un p-value piccolo indica che il risultato osservato è poco
compatibile con l'ipotesi nulla. Un p-value elevato, al contrario,
non dimostra che i gruppi siano identici, ma indica che i dati non
forniscono evidenza sufficiente per rifiutare l'ipotesi nulla.


\subsubsection{Cohen's \texorpdfstring{$d$}{d}}

La significatività statistica non fornisce una misura diretta
dell'ampiezza della differenza. Per questo motivo, per ogni feature è
stato calcolato anche Cohen's $d$:

\begin{equation}
	d
	=
	\frac{
		\overline{x}_{\mathrm{ADHD}}
		-
		\overline{x}_{\mathrm{Control}}
	}{
		s_{\mathrm{pooled}}
	}.
\end{equation}

Nel notebook la deviazione standard combinata è stata calcolata come:

\begin{equation}
	s_{\mathrm{pooled}}
	=
	\sqrt{
		\frac{
			s_{\mathrm{ADHD}}^2
			+
			s_{\mathrm{Control}}^2
		}{2}
	}.
\end{equation}

Questa formulazione assegna lo stesso peso alle varianze dei due
gruppi ed è appropriata nel presente dataset, nel quale le numerosità
sono quasi identiche, rispettivamente 61 e 60 soggetti.

Con la convenzione adottata:

\begin{equation}
	d>0
	\quad\Longrightarrow\quad
	\overline{x}_{\mathrm{ADHD}}
	>
	\overline{x}_{\mathrm{Control}},
\end{equation}

mentre:

\begin{equation}
	d<0
	\quad\Longrightarrow\quad
	\overline{x}_{\mathrm{Control}}
	>
	\overline{x}_{\mathrm{ADHD}}.
\end{equation}

Il valore assoluto $|d|$ descrive invece la grandezza della
separazione standardizzata ed è stato utilizzato per ordinare le
feature.

Mann--Whitney e Cohen's $d$ forniscono informazioni complementari:
il primo è basato sui ranghi e valuta l'evidenza statistica di una
differenza, mentre il secondo descrive la differenza tra le medie in
unità di deviazione standard. 


\subsubsection{Correzione per confronti multipli}

Il test di Mann--Whitney è stato ripetuto per tutte le 120 feature
spettrali. Applicando indipendentemente la soglia $p<0.05$ a molti
test, la probabilità di ottenere almeno alcuni risultati
apparentemente significativi per effetto del caso aumenta.

Per controllare questo problema è stata applicata la procedura
Benjamini--Hochberg per il controllo del
\textit{False Discovery Rate} (FDR) \cite{multiple}.

Indicando con $m$ il numero totale di test, i p-value vengono ordinati:

\begin{equation}
	p_{(1)}
	\leq
	p_{(2)}
	\leq
	\ldots
	\leq
	p_{(m)}.
\end{equation}

Ciascun p-value ordinato viene confrontato con:

\begin{equation}
	\frac{i}{m}q,
\end{equation}

dove $i$ è la sua posizione nell'ordinamento e $q=0.05$ è il livello
di FDR scelto.

La procedura individua il massimo indice $i^\ast$ per cui:

\begin{equation}
	p_{(i^\ast)}
	\leq
	\frac{i^\ast}{m}q.
\end{equation}

Vengono quindi rifiutate le ipotesi associate ai p-value fino a tale
posizione. Nel notebook sono stati inoltre calcolati direttamente i
p-value corretti, indicati come $p_{\mathrm{FDR}}$, e sono state
considerate significative le feature che soddisfacevano:

\begin{equation}
	p_{\mathrm{FDR}}<0.05.
\end{equation}

Il controllo FDR limita la proporzione attesa di falsi positivi tra i
risultati dichiarati significativi; non garantisce invece che ogni
singolo risultato significativo sia vero.


\subsubsection{Correlazione di Spearman}

Dopo il confronto tra i gruppi è stata analizzata la ridondanza tra le
feature mediante il coefficiente di correlazione di Spearman.

A differenza dell'analisi precedente, questa procedura non verifica se
una feature differisca tra ADHD e controlli. Valuta invece se i
soggetti che presentano valori elevati in una feature tendano a
presentare valori elevati, o bassi, anche in un'altra feature.

Spearman è definita come la correlazione di Pearson calcolata sui
ranghi:

\begin{equation}
	\rho_s
	=
	\operatorname{corr}
	\left[
	R(X),R(Y)
	\right],
\end{equation}

dove $R(X)$ e $R(Y)$ sono i ranghi delle osservazioni delle due
variabili.

Nel caso in cui non siano presenti ranghi uguali, il coefficiente può
anche essere espresso come:

\begin{equation}
	\rho_s
	=
	1
	-
	\frac{
		6\sum_{i=1}^{n}d_i^2
	}{
		n(n^2-1)
	},
\end{equation}

dove $d_i$ è la differenza tra i ranghi delle due variabili per il
soggetto $i$.

Il coefficiente assume valori compresi tra $-1$ e $1$:

\begin{itemize}
	\item $\rho_s=1$ indica un'associazione monotona positiva
	perfetta;
	
	\item $\rho_s=-1$ indica un'associazione monotona negativa
	perfetta;
	
	\item un valore vicino a zero indica l'assenza di una chiara
	associazione monotona.
\end{itemize}

Un valore nullo non esclude necessariamente la presenza di relazioni
non monotone.

Spearman è stata preferita a Pearson per l'analisi delle feature
spettrali perché è basata sui ranghi, non richiede normalità ed è meno
sensibile ai valori estremi.

Le correlazioni sono state calcolate anche separatamente nei gruppi
ADHD e Control, sia per le feature di canale sia per quelle regionali.
Sono state inoltre studiate separatamente le correlazioni tra le
potenze della stessa banda nei diversi canali.


\subsection{Classificazione}

Ogni soggetto è stato rappresentato mediante un vettore di feature e
la classe diagnostica è stata utilizzata come variabile target:

\begin{equation}
	y=
	\begin{cases}
		0, & \text{Control},\\
		1, & \text{ADHD}.
	\end{cases}
\end{equation}

Sono stati considerati due classificatori:

\begin{itemize}
	\item Random Forest, come modello capace di rappresentare
	relazioni non lineari e interazioni tra feature;
	
	\item Logistic Regression, come modello lineare di confronto.
\end{itemize}


\subsubsection{Suddivisione dei dati e cross-validation}

Prima di qualsiasi confronto tra feature set o ottimizzazione degli
iperparametri, i 121 soggetti sono stati suddivisi in:

\begin{itemize}
	\item training set, pari all'80\% del campione;
	\item test set indipendente, pari al 20\% del campione.
\end{itemize}

La divisione è stata eseguita mediante campionamento stratificato, in
modo da mantenere proporzioni simili di ADHD e controlli nei due
sottoinsiemi, utilizzando:

\begin{equation}
	\texttt{random\_state}=42.
\end{equation}

Il test set non è stato utilizzato per confrontare i feature set,
scegliere il modello o ottimizzare gli iperparametri.

Sul solo training set è stata applicata una cross-validation
stratificata a cinque fold, con mescolamento dei soggetti e
\texttt{random\_state}=42. In ciascuna iterazione, quattro fold sono
stati utilizzati per l'addestramento e il quinto per la validazione.
La procedura è stata ripetuta cinque volte, affinché ciascun fold
fosse utilizzato una volta come validation set.

Dopo la selezione della configurazione migliore, il modello è stato
addestrato nuovamente sull'intero training set e valutato una sola
volta sul test set.


\subsubsection{Metriche di valutazione}

Sono state considerate le seguenti metriche:

\begin{equation}
	\mathrm{Accuracy}
	=
	\frac{TP+TN}{TP+TN+FP+FN},
\end{equation}

\begin{equation}
	\mathrm{Recall}
	=
	\frac{TP}{TP+FN},
\end{equation}

\begin{equation}
	\mathrm{Precision}
	=
	\frac{TP}{TP+FP},
\end{equation}

\begin{equation}
	F_1
	=
	2
	\frac{
		\mathrm{Precision}\cdot\mathrm{Recall}
	}{
		\mathrm{Precision}+\mathrm{Recall}
	}.
\end{equation}

La balanced accuracy è la media del recall delle due classi:

\begin{equation}
	\mathrm{Balanced\ Accuracy}
	=
	\frac{
		\mathrm{Sensitivity}
		+
		\mathrm{Specificity}
	}{2}.
\end{equation}

È stata inoltre calcolata la ROC-AUC, che misura la capacità del
modello di assegnare punteggi più elevati ai soggetti ADHD rispetto
ai controlli considerando tutte le possibili soglie decisionali.


\subsubsection{Random Forest}

La Random Forest è un metodo ensemble costituito da numerosi alberi
decisionali.

Ogni albero viene addestrato su un campione del training
set, ottenuto mediante estrazione casuale con reinserimento. Alcuni
soggetti possono quindi comparire più volte nel campione di uno
specifico albero, mentre altri possono non essere inclusi.

Inoltre, a ogni nodo viene valutato soltanto un sottoinsieme casuale
delle feature. Questa seconda fonte di casualità riduce la
correlazione tra gli alberi: alberi differenti tendono a costruire
regole differenti e a commettere errori non perfettamente coincidenti.

La combinazione delle previsioni riduce soprattutto la varianza
rispetto a un singolo albero. Nell'implementazione di
\texttt{scikit-learn}, le probabilità di classe prodotte dai singoli
alberi vengono mediate e viene selezionata la classe con probabilità
media maggiore.

Per il confronto iniziale tra feature set è stata utilizzata una
Random Forest con:

\begin{verbatim}
	n_estimators = 500
	max_depth = None
	min_samples_leaf = 2
	class_weight = "balanced"
	random_state = 42
\end{verbatim}

Il parametro \texttt{class\_weight="balanced"} assegna pesi
inversamente proporzionali alle numerosità delle classi. Sebbene il
dataset sia quasi perfettamente bilanciato, il parametro è stato
mantenuto per rendere il modello robusto alle piccole differenze
presenti nei singoli fold.


\paragraph{Ottimizzazione degli iperparametri}

Il tuning è stato effettuato esclusivamente sul training set in due
fasi.

Nella prima fase è stata utilizzata una
\texttt{RandomizedSearchCV}, valutando 80 combinazioni estratte da uno
spazio ampio:

\begin{verbatim}
	n_estimators:       valori interi tra 200 e 1499
	max_depth:          None, 3, 5, 7, 10, 15, 20, 30
	min_samples_leaf:   valori interi tra 1 e 9
	min_samples_split:  valori interi tra 2 e 14
	max_features:       "sqrt", "log2", None
\end{verbatim}

La ricerca randomizzata permette di esplorare lo spazio degli
iperparametri senza valutare esaustivamente tutte le possibili
combinazioni.

Nella seconda fase è stata applicata una
\texttt{GridSearchCV} locale costruita intorno ai valori migliori
individuati nella prima ricerca. Questa fase ha permesso di esplorare
sistematicamente una regione più ristretta dello spazio degli
iperparametri.

Entrambe le ricerche hanno utilizzato la cross-validation
stratificata a cinque fold e la balanced accuracy come funzione
obiettivo.


\paragraph{Impurità di Gini e feature importance}

In un nodo contenente proporzioni $p_k$ delle diverse classi,
l'impurità di Gini è:

\begin{equation}
	G(v)
	=
	1-\sum_{k=1}^{K}p_{k,v}^{\,2}.
\end{equation}

Nel caso binario:

\begin{equation}
	G(v)
	=
	1
	-
	p_{\mathrm{ADHD},v}^{\,2}
	-
	p_{\mathrm{Control},v}^{\,2}.
\end{equation}

L'impurità è nulla quando tutti i soggetti nel nodo appartengono alla
stessa classe e raggiunge il massimo di 0.5 quando le due classi sono
presenti in uguale proporzione.

Per uno split che divide un nodo padre $v$ nei nodi figli $v_L$ e
$v_R$, la riduzione di impurità è:

\begin{equation}
	\Delta G(v)
	=
	G(v)
	-
	\left[
	\frac{N_L}{N_v}G(v_L)
	+
	\frac{N_R}{N_v}G(v_R)
	\right].
\end{equation}

Il contributo dello split viene pesato per la proporzione di soggetti
che raggiungono il nodo. Gli split vicini alla radice tendono quindi
ad avere un peso maggiore rispetto a quelli che coinvolgono pochi
soggetti.

L'importanza di una feature è ottenuta sommando le riduzioni pesate di
impurità associate agli split che utilizzano quella feature, mediando
il risultato su tutti gli alberi e normalizzando le importanze in modo
che la loro somma sia uguale a uno.

La feature importance non indica la direzione dell'associazione e non
deve essere interpretata come effetto causale. Inoltre, in presenza di
feature correlate, più variabili possono essere utilizzate
alternativamente nei diversi alberi e condividere la stessa
importanza.


\subsubsection{Logistic Regression}

La Logistic Regression è stata utilizzata come modello lineare di
confronto.

Il modello calcola una combinazione lineare delle feature:

\begin{equation}
	z
	=
	w_0
	+
	\sum_{j=1}^{p}w_jx_j,
\end{equation}

e la trasforma mediante la funzione sigmoide:

\begin{equation}
	P(y=1\mid\mathbf{x})
	=
	\sigma(z)
	=
	\frac{1}{1+e^{-z}}.
\end{equation}



La standardizzazione è stata inserita all'interno di una pipeline:

\begin{equation}
	z_j
	=
	\frac{x_j-\mu_j}{s_j}.
\end{equation}

Durante ogni fold, media e deviazione standard sono state calcolate
esclusivamente sui dati di training e successivamente applicate al
validation set. Questa procedura evita che l'informazione contenuta
nei dati di validazione influenzi il preprocessing.


\paragraph{Regolarizzazione e tuning}

È stata utilizzata una penalizzazione $L_2$:

\begin{equation}
	\lambda
	\sum_{j=1}^{p}w_j^2,
\end{equation}

che riduce l'ampiezza dei coefficienti e limita la complessità del
modello, senza effettuare normalmente una selezione esatta delle
feature.

In \texttt{scikit-learn}, la forza della penalizzazione è regolata
mediante il parametro:

\begin{equation}
	C
	=
	\frac{1}{\lambda}.
\end{equation}

Valori piccoli di $C$ corrispondono quindi a una regolarizzazione più
forte, mentre valori grandi corrispondono a una penalizzazione più
debole.

Il parametro è stato ottimizzato mediante Grid Search sui valori:

\begin{equation}
	C
	\in
	\{
	0.001,\,
	0.01,\,
	0.1,\,
	1,\,
	10,\,
	100
	\}.
\end{equation}

Sono stati utilizzati penalizzazione $L_2$, solver
\texttt{lbfgs}, \texttt{max\_iter}=5000 e
\texttt{class\_weight="balanced"}. La configurazione migliore è stata
selezionata mediante balanced accuracy nella cross-validation
stratificata a cinque fold.


\subsection{Feature di functional connectivity}

L'analisi spettrale descrive la quantità di potenza presente in ogni
canale, ma non descrive direttamente come due segnali varino insieme
nel tempo. Per ottenere una rappresentazione complementare è stata
quindi calcolata la functional connectivity EEG.

Nel presente lavoro, la functional connectivity è definita come
dipendenza statistica tra serie temporali registrate in canali o
regioni differenti. Essa non dimostra l'esistenza di un collegamento
anatomico diretto, non stabilisce una direzione dell'interazione e non
permette di inferire causalità.


\subsubsection{Filtraggio nelle singole bande}

A partire dal segnale già preprocessato, per ciascun soggetto sono
stati nuovamente isolati i segnali delle singole bande mediante un
filtro Butterworth passa-banda di ordine quattro:

\begin{table}[H]
	\centering
	\begin{tabular}{lc}
		\hline
		\textbf{Banda} & \textbf{Intervallo} \\
		\hline
		Delta & 1--4 Hz \\
		Theta & 4--8 Hz \\
		Alpha & 8--13 Hz \\
		Beta & 13--30 Hz \\
		\hline
	\end{tabular}
	\caption{Bande utilizzate nell'analisi di functional
		connectivity.}
	\label{tab:connectivity_bands}
\end{table}

Il filtro è stato applicato mediante \texttt{filtfilt}, eseguendo il
filtraggio prima in avanti e poi all'indietro. Il risultato presenta
un ritardo di fase netto nullo, caratteristica importante perché la
connettività viene calcolata a partire dalle relazioni temporali tra
i canali.


\subsubsection{Correlazione di Pearson}

Per due serie temporali $x(t)$ e $y(t)$, il coefficiente di Pearson è:

\begin{equation}
	r_{xy}
	=
	\frac{
		\sum_{t=1}^{T}
		\left(x_t-\overline{x}\right)
		\left(y_t-\overline{y}\right)
	}{
		\sqrt{
			\sum_{t=1}^{T}
			\left(x_t-\overline{x}\right)^2
		}
		\sqrt{
			\sum_{t=1}^{T}
			\left(y_t-\overline{y}\right)^2
		}
	}.
\end{equation}

Il coefficiente assume valori compresi tra $-1$ e $1$:

\begin{itemize}
	\item $r>0$: i due segnali tendono a variare nella stessa
	direzione;
	
	\item $r<0$: i due segnali tendono a variare in direzioni
	opposte;
	
	\item $r\simeq0$: non emerge una chiara relazione lineare a
	ritardo nullo.
\end{itemize}

Un valore prossimo a zero non esclude una relazione non lineare o una
relazione caratterizzata da un ritardo temporale. Inoltre, correlazioni
elevate possono essere influenzate dalla conduzione di volume, dalla
scelta del riferimento e da sorgenti comuni.


\subsubsection{Connettività tra canali}

Per ciascuna banda è stata calcolata la correlazione di Pearson tra
tutte le coppie uniche dei 19 canali.

Il numero di coppie è:

\begin{equation}
	\binom{19}{2}
	=
	\frac{19\cdot18}{2}
	=
	171.
\end{equation}

Poiché sono state analizzate quattro bande, il numero totale di
feature a livello di canale è:

\begin{equation}
	171\cdot4
	=
	684.
\end{equation}

Ogni feature rappresenta quindi la correlazione temporale a ritardo
nullo tra due canali in una specifica banda. Un esempio è:

\begin{center}
	\texttt{F3\_C3\_theta\_connectivity}.
\end{center}


\subsubsection{Connettività regionale}

Per costruire la rappresentazione regionale, dopo il filtraggio sono
stati mediati campione per campione i segnali dei canali appartenenti
alla stessa regione:

\begin{equation}
	x_R(t)
	=
	\frac{1}{|R|}
	\sum_{c\in R}
	x_c(t).
\end{equation}

Sono state così ottenute cinque serie temporali regionali: frontale,
centrale, temporale, parietale e occipitale.

La correlazione di Pearson è stata quindi calcolata tra ciascuna
coppia di serie regionali. Non è stata quindi calcolata la media delle
correlazioni tra tutte le coppie di elettrodi delle due regioni:
prima sono stati mediati i segnali e successivamente è stata calcolata
la correlazione tra i segnali regionali risultanti.

Con cinque regioni si ottengono:

\begin{equation}
	\binom{5}{2}
	=
	10
\end{equation}

coppie per banda e:

\begin{equation}
	10\cdot4
	=
	40
\end{equation}

feature regionali.


I  feature set sono stati confrontati mediante la stessa procedura
di classificazione descritta per le feature spettrali: split
stratificato 80/20, cross-validation stratificata a cinque fold sul
training set, selezione mediante balanced accuracy, tuning degli
iperparametri e valutazione finale sul test set indipendente.